\documentclass[a4paper,11pt]{article}
\usepackage[utf8]{inputenc}
\usepackage[T1]{fontenc}
\usepackage[french]{babel}
\usepackage[left=1cm,right=1cm,top=.5cm,bottom=0cm]{geometry}
\usepackage{enumitem}
\usepackage{hyperref}
\usepackage{xcolor}
\usepackage{titlesec}
\usepackage{fontawesome5}
\usepackage{lmodern}
\usepackage[normalem]{ulem}

% --- Couleurs Thales ---
\definecolor{primary}{RGB}{0, 45, 114}
\definecolor{secondary}{RGB}{80, 80, 80}

% --- Config Liens ---
\hypersetup{colorlinks=true, linkcolor=primary, filecolor=primary, urlcolor=primary}

% --- Titres ---
\titleformat{\section}{\large\bfseries\color{primary}\uppercase}{}{0em}{}[\titlerule]
\titlespacing{\section}{0pt}{8pt}{5pt}

\begin{document}

% --- EN-TÊTE ---
\begin{center}
    {\Large \textbf{Loïc Aron MBASSI EWOLO}} \\ \vspace{4pt}
    \textbf{\large Alternance Ingénieur Développement Logiciel Fullstack} \\
    \textit{\small Disponible dès septembre 2026 pour 2 à 3 ans} \\ \vspace{4pt}
    \small
    \faMapMarker\ Cergy, France (Mobile IDF) \quad
    \faPhone\ (+33) 7 59 52 33 98 \quad
    \faEnvelope\ \href{mailto:loicaron.mbassiewolo@ensea.fr}{loicaron.mbassiewolo@ensea.fr} \\ \vspace{2pt}
    \faLinkedin\ \href{https://www.linkedin.com/in/loic-mbassi/}{linkedin.com/in/loic-mbassi} \quad
    \faGithub\ \href{https://github.com/Nameless0l}{github.com/Nameless0l} \quad
    \faGlobe\ \href{https://mbassiloic.tech}{mbassiloic.tech}
\end{center}

\vspace{-6pt}

% --- PROFIL ---
\noindent\small{Élève-ingénieur double diplôme ENSEA/Polytechnique Yaoundé avec \textbf{+3 ans d'expérience en développement web PHP/Laravel} et pratique de Java (Spring Boot, Design Patterns). Expertise en APIs REST, bases de données SQL, architecture logicielle et intégration continue. Stage de recherche INRIA Saclay en cours (avril-août 2026). Habitué au travail en mode Agile, à l'analyse de besoins métier et à la documentation technique.}

% --- FORMATION ---
\section{Formation}
\begin{itemize}[leftmargin=*, label=\tiny$\bullet$, nosep]
    \item \textbf{ENSEA -- École Nationale Supérieure de l'Électronique et de ses Applications} \hfill \textit{2025 -- 2027} \\
    \small{Double diplôme d'Ingénieur -- Systèmes Embarqués, Signal, IA. Cursus incluant architecture logicielle et DevOps.}
    \item \textbf{École Polytechnique de Yaoundé (1\textsuperscript{ère} école d'ingénieur CEMAC)} \hfill \textit{2021 -- 2025} \\
    \small{Diplôme d'Ingénieur de Conception (Master 2) : Génie Logiciel, POO avancée, Design Patterns, Algorithmique, BDD.}
\end{itemize}

% --- COMPÉTENCES ---
\section{Compétences Techniques}
\begin{itemize}[leftmargin=*, nosep]
    \item \small{\textbf{Backend :} PHP 8 (maîtrise), Laravel (expert), Symfony (notions), Java (Spring Boot), APIs RESTful, OAuth2/JWT.}
    \item \small{\textbf{Frontend :} HTML5, CSS3/SASS, JavaScript/TypeScript, React.js (maîtrise), Responsive Design.}
    \item \small{\textbf{Bases de données :} MySQL (maîtrise), SQL avancé, MongoDB, Redis. Modélisation relationnelle et optimisation.}
    \item \small{\textbf{DevOps \& Outils :} Git/GitHub/GitLab (maîtrise), Docker, CI/CD, Tests unitaires (PHPUnit), Jira, Documentation.}
    \item \small{\textbf{Protocoles \& Sécurité :} HTTP/HTTPS, SSH, FTP, LDAP (notions), Sécurisation d'APIs, RGPD.}
    \item \small{\textbf{Méthodologies :} Agile/Scrum/Kanban, Clean Code, SOLID, Code review, Rédaction de spécifications.}
    \item \small{\textbf{Langues :} Français (natif), Anglais (professionnel/technique -- collaboration internationale).}
\end{itemize}

% --- EXPÉRIENCES PROFESSIONNELLES ---
\section{Expériences Professionnelles}
\begin{itemize}[leftmargin=*, label={-}]

\item \textbf{Stagiaire Recherche @ INRIA Saclay -- Équipe TRIBE} \hfill \textit{Avril 2026 -- Août 2026} \\
\textit{Palaiseau (Plateau de Saclay) -- Sécurité numérique / Privacy} \hfill \textit{Python, NLP, pandas, Git}
\vspace{-2pt}
    \begin{itemize}[leftmargin=0.4cm, nosep, label=\tiny$\bullet$]
        \item \small{Conception d'un pipeline Python d'analyse automatisée d'applications mobiles (permissions, SDKs, politiques de confidentialité).}
        \item \small{Traitement de données à grande échelle et rédaction scientifique en anglais. Collaboration multi-laboratoires (INRIA/LVMT).}
    \end{itemize}
\vspace{3pt}

\item \textbf{Développeur Backend @ CSC (Cameroon Software Company)} \hfill \textit{Oct 2024 -- Jan 2025} \\
\textit{Yaoundé -- Secteur Fintech / Bancaire} \hfill \textit{Laravel, MySQL, Docker, REST APIs, Redis}
\vspace{-2pt}
    \begin{itemize}[leftmargin=0.4cm, nosep, label=\tiny$\bullet$]
        \item \small{Conception et développement du backend d'une plateforme de transactions financières sécurisée.}
        \item \small{Implémentation d'APIs RESTful documentées pour les intégrations tierces.}
        \item \small{Gestion de transactions concurrentes et conformité aux normes de sécurité financière.}
        \item \small{Travail en équipe Agile : sprints, démonstrations, revues de code, documentation technique.}
    \end{itemize}
\vspace{3pt}

\item \textbf{Développeur Web Full Stack @ SOSDOCTA} \hfill \textit{Fév 2022 -- Mars 2024} \\
\textit{Remote (Belgique/Canada) -- Télémédecine} \hfill \textit{Laravel, PHP, MySQL, JS, OAuth2/JWT}
\vspace{-2pt}
    \begin{itemize}[leftmargin=0.4cm, nosep, label=\tiny$\bullet$]
        \item \small{Développement complet d'une plateforme de télémédecine : backend, dashboards, système de paiement.}
        \item \small{Maintenance évolutive sur 2 ans : corrections, nouvelles fonctionnalités, sécurité.}
        \item \small{Gestion sécurisée des données médicales sensibles (conformité RGPD, OAuth2/JWT).}
    \end{itemize}
\vspace{3pt}

\item \textbf{Développeur Full Stack @ NettoieQuebec \& BAZEL SQUARE} \hfill \textit{2023 -- 2025} \\
\textit{Remote (Québec) / Yaoundé -- Services / E-commerce} \hfill \textit{Laravel, PHP, MySQL, JavaScript}
\vspace{-2pt}
    \begin{itemize}[leftmargin=0.4cm, nosep, label=\tiny$\bullet$]
        \item \small{Développement from scratch de plateformes web avec parcours utilisateur complets et paiements intégrés.}
        \item \small{Conception de systèmes de réservation, gestion de catalogue et notifications automatisées.}
    \end{itemize}
\vspace{3pt}

\end{itemize}

% --- PROJETS TECHNIQUES ---
\section{Projets Techniques}
\begin{itemize}[leftmargin=*, label={-}]

\item \textbf{Laravel API Generator (Open Source -- +30 \faGithub\ stars)} \hfill \textit{2024 -- Présent} \\
\textit{Projet personnel publié sur Packagist} \hfill \textit{PHP, Laravel, Python, XML, LLM}
\vspace{-2pt}
    \begin{itemize}[leftmargin=0.4cm, nosep, label=\tiny$\bullet$]
        \item \small{Outil de génération automatique de code backend Laravel (modèles, contrôleurs, migrations, routes) depuis UML/XML.}
        \item \small{Parsing géométrique avancé et intégration LLM. +20 téléchargements sur Packagist.}
    \end{itemize}
\vspace{3pt}

\item \textbf{GoFind -- Architecture Microservices} \hfill \textit{2024} \\
\textit{Projet académique} \hfill \textit{Laravel, NestJS, MongoDB, RabbitMQ, Docker}
\vspace{-2pt}
    \begin{itemize}[leftmargin=0.4cm, nosep, label=\tiny$\bullet$]
        \item \small{Conception d'une architecture microservices complète avec communication asynchrone (RabbitMQ) et API Gateway.}
        \item \small{Containerisation Docker et intégration continue pour le déploiement automatisé.}
    \end{itemize}
\vspace{3pt}

\item \textbf{Freengo -- Application Covoiturage (Uber-like)} \hfill \textit{2024} \\
\textit{Projet académique} \hfill \textit{Spring Boot (Java), ScyllaDB, Next.js}
\vspace{-2pt}
    \begin{itemize}[leftmargin=0.4cm, nosep, label=\tiny$\bullet$]
        \item \small{Backend Java Spring Boot avec base NoSQL haute performance et algorithmes d'optimisation de routes.}
    \end{itemize}
\vspace{3pt}

\end{itemize}

% --- DISTINCTIONS ---
\section{Distinctions et Leadership}
\begin{itemize}[leftmargin=*, label=\tiny$\bullet$, nosep]
    \item \small{\textbf{1\textsuperscript{ère} Place} -- IndabaX Africa 2025 (Hackathon IA Santé) : déploiement API Flask, dashboard Streamlit.}
    \item \small{\textbf{1\textsuperscript{ère} Place} -- CITS National Hackathon 2024 (75 équipes) : optimisation collecte déchets urbains.}
    \item \small{\textbf{Lead AI Cell} @ Polytechnique Yaoundé : +50\% de membres, workshops Laravel/React/ML.}
\end{itemize}

\end{document}
